\begin{abstract}
Federated learning is promoted for household load forecasting because it keeps smart-meter data at home. Yet its
accuracy loss is usually blamed on data heterogeneity, and its bandwidth advantage is assumed rather than measured.
We test both claims on 100 London households from the Low Carbon London trial with a controlled protocol: the same
homes, test set, LSTM and compute budget, changing one factor at a time: number of clients, heterogeneity (dialled
over real load-profile clusters and measured by a label-free index), participation, aggregation rule (FedAvg, FedProx,
a heterogeneity-adaptive proximal rule) and update precision.
Every byte exchanged is counted. Federated training on a single client matches centralised accuracy. The gap is
caused by fragmentation, growing from 2\,\% with five clients to 10\,\% with one client per home. Heterogeneity at
a fixed number of clients barely changes final error, but raises the traffic needed to approach centralised accuracy
by about 85\,\%. Federated training moves more bytes than uploading the raw readings. FedAvg leaves most households
worse off than local models; one epoch of local fine-tuning largely removes this at community scale. Proximal
regularisation does not help, and its adaptive variant fails because update dissimilarity is dominated by convergence,
not heterogeneity.
\end{abstract}
